\label{chap:log-ht-monodromy-frontier}

\providecommand{\A}{\mathcal A}
\providecommand{\B}{\mathcal B}
\providecommand{\C}{\mathcal C}
\providecommand{\D}{\mathcal D}
\providecommand{\E}{\mathcal E}
\providecommand{\F}{\mathcal F}
\providecommand{\FM}{\mathrm{FM}}
\providecommand{\Conf}{\mathrm{Conf}}
\providecommand{\HT}{\mathrm{HT}}
\providecommand{\MC}{\mathrm{MC}}
\providecommand{\End}{\mathrm{End}}
\providecommand{\Hom}{\mathrm{Hom}}
\providecommand{\id}{\mathrm{id}}
\providecommand{\im}{\mathrm{im}}
\providecommand{\Res}{\mathrm{Res}}
\providecommand{\Aut}{\mathrm{Aut}}
\providecommand{\PP}{\mathbb P}
\providecommand{\AA}{\mathbb A}
\providecommand{\CC}{\mathbb C}
\providecommand{\RR}{\mathbb R}
\providecommand{\ZZ}{\mathbb Z}
\providecommand{\kk}{\Bbbk}
\providecommand{\h}{\hbar}
\providecommand{\ot}{\otimes}
\providecommand{\wot}{\widehat\otimes}
\providecommand{\s}{\mathrm s}
\providecommand{\ds}{d_{\mathrm{dR}}}
\providecommand{\CYB}{\mathrm{CYB}}
\providecommand{\PT}{\mathrm{PT}}
\providecommand{\Disk}{\mathrm{Disk}}
\providecommand{\SN}{\boldsymbol{\nabla}}
\providecommand{\bboxprod}{\mathop{\boxtimes}\displaylimits}

\title{Logarithmic HT Monodromy: Extensions and Frontier}
\author{}
\date{}

\section{Conjectural bridges}
The theorem package above strongly suggests further identifications.

\begin{conjecture}[Monodromy equals spectral $R$-matrix; \ClaimStatusConjectured]\label{conj:rmatrix}
The monodromy of the reduced HT logarithmic connection is the spectral $R$-matrix controlling the factorisation quantum group attached to the same dg-shifted Yangian data. Equivalently, the monodromy functor of the reduced connection should identify the reduced HT line category with the representation category of the associated factorisation quantum group.
\end{conjecture}

\begin{conjecture}[Logarithmic CoHA envelope; \ClaimStatusConjectured]\label{conj:coha}
For a symmetric quiver there exists a logarithmic factorization enhancement of the cohomological Hall algebra on logarithmic Fulton-MacPherson spaces whose connected collision residues recover the primitive Lie algebra of the associated vertex algebra, and whose Koszul dual is the dg-shifted Yangian governing the corresponding line-operator theory.
\end{conjecture}

\begin{conjecture}[Chromatic lift of the HT associator; \ClaimStatusConjectured]\label{conj:chromatic-lift}
For a spectral lift of a resolved logarithmic HT datum, the reduced HT associator admits compatible chromatic localizations whose height-$n$ pieces are computed by the monodromy of the $E(n)$-localized spectral factorization algebra.
\end{conjecture}


\section{The Kazhdan--Lusztig argument for the monodromy-braiding identification}
\label{sec:KL-monodromy-braiding}

\providecommand{\fS}{\mathfrak{S}}
\providecommand{\cL}{\mathcal{L}}
\providecommand{\cM}{\mathcal{M}}
\providecommand{\Nil}{\widetilde{\mathcal{N}}}
\providecommand{\fg}{\mathfrak{g}}

In geometric representation theory, the identification of KZ monodromy with the $R$-matrix of the quantum group is the Kazhdan--Lusztig theorem \cite{kazhdan-lusztig-iv}.
Their proof proceeds by computing the monodromy of the KZ connection via convolution on the Steinberg variety $Z = \Nil \times_{\fg} \Nil$: the flat connection on category $\mathcal{O}$ is realized geometrically as the Gauss--Manin connection on the family of Steinberg fibers, and its monodromy is read off from the convolution product, which is the quantum group multiplication.
The chiral bar complex gives us the same ingredients in the HT setting.
The Steinberg object $\fS_b = \cL_b \times_{\cM_{\mathrm{vac}}} \cL_b$ (the self-fiber-product of the boundary line category over the vacuum moduli) plays the role of $Z$; the $(-1)$-shifted symplectic structure plays the role of the moment map; and the bar differential encodes the convolution.
We now make this precise.

\begin{construction}[The flat connection on $\fS_b$]
\label{constr:steinberg-connection}
The $(-1)$-shifted symplectic form $\omega_{-1}$ on the Steinberg object $\fS_b = \cL_b \times_{\cM_{\mathrm{vac}}} \cL_b$ carries a canonical Liouville $1$-form $\theta$ (the primitive of the shifted symplectic structure, existing because the fiber product is Lagrangian in the appropriate sense).
When the HT theory is defined on a punctured disk $D^* \subset \CC$, the family $\fS_b \to D^*$ acquires a flat connection
\[
\nabla_{\fS} \;=\; d \;+\; r(z)\,\frac{dz}{z}
\]
where $r(z)$ is the classical $r$-matrix: the residue of the chiral OPE along the collision divisor $\Delta \subset D \times D$.
Concretely, $r(z)$ is the element of $\End(\cL_b^{\ot 2})$ whose matrix coefficients are the polar parts of the OPE coefficients $\langle a(w)\, b(z) \rangle \sim \sum_{n \geq 0} (a_{(n)} b)(z) / (w-z)^{n+1}$.
The flatness $\nabla_{\fS}^2 = 0$ is equivalent to the classical Yang--Baxter equation $[\CYB](r) = 0$, which holds because $r(z)$ is the structure constant of a vertex algebra OPE (the Jacobi identity implies CYB).

The connection is regular singular at $z = 0$: the pole is simple because the OPE has at worst first-order poles after extracting the leading singularity.
The monodromy around $z = 0$ is therefore
\[
\mathrm{Mon}_{z=0}(\nabla_{\fS}) \;=\; \exp\bigl(2\pi i \cdot \Res_{z=0}\, r(z)\,dz/z\bigr) \;=\; \exp(2\pi i\, r_0)
\]
where $r_0 = \Res_{z=0} r(z)\,dz/z$ is the residue endomorphism.
\end{construction}

\begin{proposition}[Free $\beta\gamma$ verification; \ClaimStatusProvedHere]
\label{prop:betagamma-monodromy}
For the free $\beta\gamma$ system on $\CC_z \times \RR_t$, the monodromy of $\nabla_{\fS}$ around $z = 0$ equals the spectral $R$-matrix $R(z) = 1 + r/z + O(1/z^2)$ of the Heisenberg vertex algebra.
\end{proposition}

\begin{proof}
The $\beta\gamma$ OPE is $\beta(w)\gamma(z) \sim 1/(w-z)$, so $r(z) = r_0/z$ with $r_0 = \beta \ot \gamma - \gamma \ot \beta \in \End(V^{\ot 2})$ the standard Heisenberg $r$-matrix.
The connection $\nabla_{\fS} = d + (r_0/z)\,dz/z$ on $D^*$ is regular singular with residue $r_0$.
Its monodromy is
\[
\mathrm{Mon}_{z=0}(\nabla_{\fS}) \;=\; \exp(2\pi i\, r_0).
\]
On the other hand, the spectral $R$-matrix of the rank-one Heisenberg algebra is $R(z) = \exp(r_0/z)$ evaluated at $z = 1$ after analytic continuation, and its monodromy representation on $\mathrm{Rep}(\mathrm{Heis})$ is precisely $\exp(2\pi i\, r_0)$.
The two expressions coincide.
\end{proof}

\begin{strategy}[Reduction to analytic input; \ClaimStatusConjectured]
\label{strat:general-monodromy}
For a general HT theory satisfying (H1)--(H4), the identification $\mathrm{Mon}_{z=0}(\nabla_{\fS}) = R(z)$ reduces to showing two things:
\begin{enumerate}
\item The connection $\nabla_{\fS}$ is regular singular at $z = 0$, with residue given by the Laplace transform of the $\lambda$-bracket:
\[
r_0 \;=\; \Res_{z=0}\, r(z)\,\frac{dz}{z}, \qquad r(z) \;=\; \int_0^\infty e^{-\lambda z}\, \{{\cdot}\,_\lambda\,{\cdot}\}\, d\lambda,
\]
which is the DK-0 formula proved in the core chapter (\S\ref{sec:DK-0}).
\item The exponential $\exp(2\pi i\, r_0)$ converges in the completed tensor product $\End(\cL_b^{\wot 2})$ and defines a braided monoidal structure on $\cL_b\text{-}\mathrm{mod}$.
\end{enumerate}
The first point is the content of the DK-0 bridge: the classical $r$-matrix IS the Laplace transform of the PVA $\lambda$-bracket, and this is proved.
The second point is an analytic convergence statement that holds whenever the $\lambda$-bracket has at most polynomial growth in $\lambda$ --- a condition guaranteed by (H2) (exponential decay of the propagator in $\RR$).
Granting this, the monodromy computation follows from the same matrix exponential as in the free field case.
\end{strategy}

\begin{remark}[Geometric origin of the $(2,0)$-projection]
\label{rmk:20-projection-origin}
The Steinberg presentation theorem of the core chapter identifies the line category $\cL_b$ with modules over the convolution algebra $H_*(\fS_b)$, and the algebraic projection $\alpha_T^{(2,0)}$ --- the $(2,0)$-component of the HT associator --- with the class of the diagonal $[\Delta] \in H_*(\fS_b)$.
Proposition~\ref{prop:betagamma-monodromy} gives this projection its geometric meaning: $\alpha_T^{(2,0)}$ is the \emph{monodromy} of the Steinberg connection $\nabla_{\fS}$, computed as the exponential of the residue $r$-matrix.
The algebraic projection and the analytic monodromy are two descriptions of the same invariant, related by the passage from the Steinberg variety to its family over $D^*$.
This is the chiral analog of the fact that, in the classical Kazhdan--Lusztig theorem, the quantum group $R$-matrix is simultaneously an algebraic object (convolution on $Z$) and an analytic object (monodromy of KZ).
\end{remark}


\section{Conclusion}\label{sec:log-ht-conclusion}

The theory developed here may be read as a chain of implications.
\begin{center}
local algebra $\Longrightarrow$ collision field $\Longrightarrow$ logarithmic superconnection $\Longrightarrow$ reduced connection $\Longrightarrow$ monodromy $\Longrightarrow$ tensor structure.
\end{center}

At the strict end, one sees a shifted KZ/FM connection. At the derived end, one sees a bar-valued superconnection whose curvature is Maurer-Cartan. Analytically, one sees renormalized graph forms on mixed logarithmic configuration spaces and compactified Stokes identities. After reduction, one sees an ordinary logarithmic connection, and from it an associator, a braiding, and the geometry of tangential basepoints. Spectrally, one sees periodicity and chromaticity become compatible with factorization homology. Across the whole theory, the true invariant object is not merely an algebra or a category, but a monodromic local-to-global machine organized by logarithmic collision geometry.

The constructions are deliberately first-principles and modular. Each piece can be strengthened independently: the mixed log-BPHZ analysis can be made fully analytic in wider classes of theories; the resolved-line criterion can be refined and classified; the monodromy/factorisation-quantum-group bridge can be proved; the logarithmic CoHA envelope can be built; and the spectral lift can be made genuinely chromatic. But the spine is now visible and, more importantly, it is rigid: once the classical line algebra resolves cleanly, the rest of the structure is no longer optional.

